\BourbakiExerciseGroup

\begin{enumerate}
\def\labelenumi{\arabic{enumi})}
\tightlist
\item
  设 \(\mathbf{K}\) 是特征为 \(p > 0\) 的域，E 是 \(\mathbf{K}\) 的一个可分扩张，且在 \(\mathbf{K}\) 上具有有限超越次数。
\end{enumerate}

\begin{enumerate}
\def\labelenumi{\alph{enumi})}
\item
  若存在 E 在 K 上的一个超越基 \(B_{0}\) 以及一个整数 \(m \geqslant 0\) ，使得 \(\mathrm{K}(\mathrm{E}^{p^{m}})\) 在 \(\mathrm{K}(\mathrm{B}_{0})\) 上可分，证明：对 E 在 K 上的每个超越基 B，存在整数 \(n \geqslant 0\) 使得 \(\mathrm{K}(\mathrm{E}^{p^{n}})\) 在 \(\mathrm{K}(\mathrm{B})\) 上可分。
\item
  由此推出：若 a) 中的条件成立，则 E 在 K 上具有一个可分超越基（若 S 是 E 在 \(\mathrm{K}(\mathrm{E}^{p})\) 上的一个 p-基，而 B 是 E 在 K 上的包含 S 的一个超越基（V，第 176 页，习题 6），则利用 V，第 176 页，习题 6，证明 B 是 E 在 K 上的一个可分超越基）。
\item
  设 K 是特征为 p \textgreater{} 0 的域，x 是 K 上的一个超越元（在 K 的一个代数闭扩张 \(\Omega\) 中）。证明诸扩张 \(\mathrm{K}(x^{p^{-n}})\) 的并 E 是 K 的一个可分扩张，满足 tr. \(\deg_{K}E = 1\) ，但不具有可分超越基。
\item
  设 K 是一个完全域。证明有理分式域 \(\mathrm{K}(\mathrm{T})\) 的完全闭包是 K 的一个超越次数为 1 的扩张，它不具有可分超越基。
\end{enumerate}

\begin{enumerate}
\def\labelenumi{\arabic{enumi})}
\setcounter{enumi}{1}
\item
  设 K 是特征为 p \textgreater{} 0 的域，E 是 K 的一个有限超越次数的扩张。要使 E 在 K 上具有一个可分超越基，必须而且只需 E 在 K 上可分，且 E 在 K 上的不完全度等于它在 K 上的超越次数。
\item
  设 \(\mathbf{E}\) 是域 \(\mathbf{K}\) 的一个扩张，\(\mathbf{F}\) 是 \(\mathbf{E}\) 的一个扩张。
\end{enumerate}

\begin{enumerate}
\def\labelenumi{\alph{enumi})}
\item
  若 E 在 K 上具有一个可分超越基，且 F 在 E 上具有一个可分超越基，则 F 在 K 上具有一个可分超越基。
\item
  若 F 在 K 上具有一个可分超越基，且 tr. \(\deg_{K}E<+\infty\) ，则 E 在 K 上具有一个可分超越基（归结为 tr. \(\deg_{K}F<+\infty\) 的情形，并应用习题 1）。
\item
  若 F 在 K 上具有一个有限可分超越基，且 F 在 E 上可分，则 F 在 E 上具有一个可分超越基（利用习题 1）。
\end{enumerate}

\begin{enumerate}
\def\labelenumi{\arabic{enumi})}
\setcounter{enumi}{3}
\item
  设 K 是特征为 p \textgreater{} 0 的域，其绝对不完全度（V，第 170 页，习题 1）等于 1。证明：要使 K 的一个扩张 E 在 K 上可分，必须而且只需 E 不包含任何在 K 上 p-根式且不属于 K 的元素（若 \(a \in K\) 构成 K 在 \(K^{p}\) 上的一个 p-基，证明 \(\mathbf{K}^{p^{-1}} = \mathbf{K}(a^{p^{-1}})\) 与 E 在 K 上线性不相交）。
\item
  \begin{enumerate}
  \def\labelenumii{\alph{enumii})}
  \tightlist
  \item
    设 K 是特征为 p \textgreater{} 0 的域，E 是 K 的一个具有可分超越基的扩张。证明诸域 \(\mathrm{K}(E^{p^{n}})\) （n 遍历 N）的交 L 是 K 的一个代数扩张。（首先考虑 \(\operatorname{tr} \cdot \deg_{K} E < +\infty\) 的情形；证明 L 在 K 上相对完全，从而 E 在 L 上可分。然后利用 V，第 176 页，习题 2 和习题 6。在一般情形，若 B 是 E 在 K 上的一个可分超越基，则对每个 \(x \in L\) ，存在 B 的一个有限子集 \(B_{0}\) ，使得 x 在每个 \(\mathrm{K}(B_{0}^{p^{n}})\) 上的极小多项式都与它在 \(\mathrm{K(B)}\) 上的极小多项式相同；由此推出，若 \(\mathrm{F} = \mathrm{K}(B_{0})(x)\) ，则对所有 n 有 \(x \in \mathrm{K}(F^{p^{n}})\) 。）
  \end{enumerate}
\end{enumerate}

\begin{enumerate}
\def\labelenumi{\alph{enumi})}
\setcounter{enumi}{1}
\tightlist
\item
  证明包含在 E 中的最大完全域 \(E^{p^{\infty}}\) （诸 \(E^{p^{n}}\) 的交）是 \(K^{p^{\infty}}\) 的一个代数扩张（利用 a) 以及 V，第 176 页，习题 8）。给出一个例子：E 是完全域 K 的有限生成扩张，但 \(E^{p^{\infty}}\) 不包含在 K 中。
\end{enumerate}

\begin{enumerate}
\def\labelenumi{\arabic{enumi})}
\setcounter{enumi}{5}
\tightlist
\item
  设 \(\mathbf{K}\) 是特征为 \(p > 0\) 的域，\(\mathbf{E}\) 是 \(\mathbf{K}\) 的一个扩张，使得存在整数 \(h \geqslant 0\) ，\(\mathbf{K}(\mathrm{E}^{p^h})\) 在 \(\mathbf{K}\) 上可分（当 \(\mathbf{E}\) 是 \(\mathbf{K}\) 的有限生成扩张时，情形总是如此）。
\end{enumerate}

\begin{enumerate}
\def\labelenumi{\alph{enumi})}
\item
  若 C 是一个 \(p\) -基（E 在 \(\mathbf{K}(\mathbf{E}^p)\) 上的），则存在 C 的一个子集 B，使得 \(\mathbf{B}^{p^h}\) 是一个 \(p\) -基（\(\mathbf{K}(\mathbf{E}^{p^h})\) 在 \(\mathbf{K}(\mathbf{E}^{p^{h+1}})\) 上的）；B 是 K 上的一个代数自由集。证明扩张 \(\mathbf{F} = \mathbf{K}(\mathbf{E}^{p^h})(\mathbf{B})\) 在 K 上可分，且 E 是 F 的代数扩张。
\item
  证明 B 是一个 \(p\) -基（F 在 \(\mathbf{K}(\mathbf{F}^p)\) 上的）。证明：若 \(x \in \mathbf{E}\) ，且 \(m\) 是使 \(x^{p^m} \in \mathbf{F}\) 的最小整数，则 \(x^{p^m} \in \mathbf{K}(\mathbf{F}^{p^m})\) ；因此 E 是域 \(\mathbf{K}^{p^{-\infty}}(\mathbf{F})\) （同构于 \(\mathbf{K}^{p^{-\infty}} \otimes_{\mathbf{K}} \mathbf{K}\) ）的一个子域。称 E 的一个子扩张 F 为可区分的，如果它在 K 上可分，且 E 是 \(\mathbf{K}^{p^{-\infty}}(\mathbf{F})\) 的一个子域；此时 F 是 K 在 E 中的一个极大可分扩张。
\item
  设 P 是特征为 p \textgreater{} 0 的完全域，\(K = P(X, Y)\) 是两个未定元的有理分式域，而 \(\Omega\) 是有理分式域 \(\mathbf{K}(\mathbf{Z})\) 在 K 上的一个代数闭扩张。若 \(u \in \Omega\) 满足 \(u^{p} = X + YZ^{p}\) ，证明 \(\mathrm{E} = \mathrm{K}(\mathrm{Z}, u)\) 是 K 的一个扩张，使得 \(\mathrm{K}(\mathrm{Z})\) 与 \(\mathrm{K}(u)\) 是 E 的两个可区分子扩张，从而 E 不存在更大的可分子扩张。若 \(v \in \Omega\) 满足 \(v^{p^{2}} = X + YZ^{p}\) ，证明在扩张 \(\mathrm{E}' = \mathrm{K}(\mathrm{Z}, v)\) 中，\(\mathrm{K}(\mathrm{Z})\) 是 E 的一个极大可分子扩张，但不是可区分子扩张。
\end{enumerate}

\begin{enumerate}
\def\labelenumi{\arabic{enumi})}
\setcounter{enumi}{6}
\tightlist
\item
  假设习题 6 的假设成立，并且 E 在 K 上的不完全度（V，第 170 页，习题 1）有限（若 E 是 K 的有限生成扩张，情形总是如此）。
\end{enumerate}

\begin{enumerate}
\def\labelenumi{\alph{enumi})}
\tightlist
\item
  设 B 是一个 \(p\) -基（F 在 \(\mathbf{K}(\mathbf{F}^p)\) 上的），它同时是 F 在 K 上的一个可分超越基（并且是 E 在 K 上的一个超越基）。证明：对每个整数 \(k > 0\) ，\(\mathbf{B}^{p^k}\) 是一个 \(p\) -无关子集（在 \(\mathbf{K}(\mathbf{E}^{p^k})\) 中、于 \(\mathbf{K}(\mathbf{E}^{p^{k+1}})\) 上）。设 \(q = \operatorname{tr} \cdot \deg_{\mathbf{K}} \mathbf{E}\) ；若 \(m_k + q\) 是 \(\mathbf{K}(\mathbf{E}^{p^k})\) 在 \(\mathbf{K}(\mathbf{E}^{p^{k+1}})\) 上的不完全度，证明 \(p^f\) —— 其中 \(f = \sum_{k=0}^{h-1} m_k\) —— 等于次数 {[}E : F{]} ，因此与
\end{enumerate}

E 的可区分子扩张 F 的选择无关。整数 \(p^{f}\) 称为 E 在 K 上的不可分度，记作 \([E:K]_{i}\) ；当 E 为有限次时，这个数与 V，第 46 页中记作 \([E:K]_{i}\) 的不可分次数一致。

\begin{enumerate}
\def\labelenumi{\alph{enumi})}
\setcounter{enumi}{1}
\item
  对 \(0 \leqslant k \leqslant h - 1\) ，\([\mathbf{K}^{p^{-k}}(\mathbf{E}) : \mathbf{K}^{p^{-k}}]_i\) 等于 \(p^{f_k}\) ，其中 \(f_k = \sum_{j=k}^{h-1} m_j\) 。
\item
  设 \(L\) 是 \(E\) 的一个子扩张，在 \(K\) 上可分，且使得 \(E\) 是 \(p\) -根式的且在 \(L\) 上有限次；证明 \([E:L] \geqslant [E:K]_i\) （若 \([L(E^{p^k}):K(E^{p^k})] = p^{r_k}\) ，证明 \(r_{k+1} \leqslant r_k + q\) ，这通过考虑 \(L(E^{p^k})\) 在 \(K(E^{p^k})\) 上的一组基以及 \(L\) 在 \(K(L^p)\) 上的一组基来完成）。
\item
  设 \(K'\) 是 K 的一个扩张，使得 \(K'\) 与 E 包含在 K 的同一个代数闭扩张 \(\Omega\) 中。证明：若 \(\mathrm{K}'(\mathrm{F})\) 在 \(K'\) 上可分，则它是 \(\mathrm{K}'(\mathrm{E})\) （ \(K'\) 的扩张）的一个可区分子扩张。当 \(K'\) 与 E 在 K 上代数不相交时即是如此，且此时我们有 \([\mathrm{K}'(\mathrm{E}):\mathrm{K}']_{i}\leqslant[\mathrm{E}:\mathrm{K}]_{i}\) ；当 \(K'\) 与 E 在 K 上线性不相交，或当 \(K'\) 是 K 的可分代数扩张时，等号成立。e) 证明 \([\mathrm{E}:\mathrm{K}]_{i}\) 是对 E 在 K 上的所有超越基 B 取的 \([\mathrm{E}:\mathrm{K}(\mathrm{B})]_{i}\) 的最小值；若 \(\bar{K}\) 是 K 在 \(\Omega\) 中的相对代数闭包（它是 K 的一个代数闭包），则 \([\mathrm{E}:\mathrm{K}(\mathrm{B})]_{i}=[\mathrm{E}:\mathrm{K}]_{i}[\bar{\mathrm{K}}(\mathrm{E}):\bar{\mathrm{K}}(\mathrm{B})]_{i}\) 。若 L 是 E 的一个子扩张，使得 E 是 L 的代数扩张，且 \(\bar{\mathrm{K}}(\mathrm{E})\) 是 \(\bar{\mathrm{K}}(\mathrm{L})\) 的（代数）可分扩张，则 L 在 E 中包含的最大可分扩张 F 是可区分的。
\item
  假设 E 是 K 的一个有限生成扩张；证明：对 E 的每个子扩张 L，我们有 \([L:K]_{i} \leqslant [E:K]_{i} \leqslant [L:K]_{i}[E:L]_{i}\) 。给出一个例子，使得
\end{enumerate}

\[
[ \mathrm{L}: \mathrm{K} ] _ {i} = [ \mathrm{E}: \mathrm{K} ] _ {i} \quad \text { and } \quad [ \mathrm{E}: \mathrm{L} ] _ {i} > 1.
\]

\begin{enumerate}
\def\labelenumi{\arabic{enumi})}
\setcounter{enumi}{7}
\tightlist
\item
  设 \(k\) 是一个域，\(L\) 是 \(k\) 的一个扩张。证明下列条件是等价的：
\end{enumerate}

\begin{enumerate}
\def\labelenumi{(\roman{enumi})}
\item
  L 是一个有限生成的纯超越扩张的可分代数扩张。
\item
  \(L\) 是 \(k\) 的一个可分扩张，且
\end{enumerate}

\[
[ \Omega_ {L / k}: L ] = \operatorname{tr} \cdot \deg_ {k} L <   + \infty .
\]

（为证明 (ii) \(\Rightarrow\) (i)，我们可以如下进行：设 L 中的 \(x_{1},\ldots ,x_{n}\) 使得 \(dx_{1},\ldots ,dx_{n}\) 构成 \(\Omega_k(\mathbf{L})\) 的一组基，并令 \(\mathbf{K} = k(x_1,\dots,x_n)\) 。首先证明 \(\operatorname {tr.deg}_k(\mathbf{K}) = n\) ，从而 K 是 \(k\) 的一个纯超越扩张。设 \(\mathbf{L}'\) 是包含在 L 中的 K 的一个有限代数扩张。通过将典范映射 \(u:\Omega_k(\mathbf{K})\otimes_\mathbf{K}\mathbf{L}'\to \Omega_k(\mathbf{L}')\) 与典范映射 \(\Omega_k(\mathbf{K})\otimes_\mathbf{K}\mathbf{L}\to \Omega_k(\mathbf{L})\) 相比较，证明 \(u\) 是单射。计算 \([\Omega_k(\mathbf{L}'):\mathbf{L}']\) （V，第 134 页，推论 1）。由此推出 \(u\) 是双射，从而 \(\Omega_{\mathbf{K}}(\mathbf{L}') = 0\) 。证明 \(\mathbf{L}'\) 在 K 上可分；由此推出 L 在 K 上可分。）

\begin{enumerate}
\def\labelenumi{\arabic{enumi})}
\setcounter{enumi}{8}
\tightlist
\item
  沿用习题 8 的记号，设 \(\mathbf{M} \subset \mathbf{L}\) 是一个子扩张。令 \(m = \operatorname{tr} \cdot \deg_k \mathbf{M}\) ，\(n = \operatorname{tr} \cdot \deg_k \mathbf{L}\) ，\(r = n - m\) 。设 \(x_1, \ldots, x_n\) （在 \(\mathbf{L}\) 中）使得 \(\mathbf{L}\) 是 \(k(x_1, \ldots, x_n) = \mathbf{K}\) 的一个可分代数扩张。
\end{enumerate}

\begin{enumerate}
\def\labelenumi{\alph{enumi})}
\item
  证明：若有必要，在将 \(x_{i}\) 重新编号后，我们可以假设 \(L\) 在 \(M(x_{1},\ldots ,x_{r})\) 上是代数的。
\item
  证明存在一个有限生成的子扩张 \(\mathbf{N} \subset \mathbf{M}\) ，使得：
\end{enumerate}

\(\alpha)\) L 在 \(\mathbf{N}(x_1,\dots,x_s)\) 上是代数的；

\(\beta)\) \(\mathbf{K}' = \mathbf{K}(\mathbf{N}(x_1,\dots,x_r))\) 是 \(\mathbf{N}(x_1,\dots,x_r)\) 的一个有限扩张；

\(\gamma)\) 典范同态 \(\mathbf{K}^{\prime}\otimes_{\mathbf{N}(x_1,\ldots ,x_r)}\mathbf{M}(x_1,\dots,x_r)\to \mathbf{K}(\mathbf{M}(x_1,\dots,x_r)) = \mathbf{K}''\) 是双射的。

\begin{enumerate}
\def\labelenumi{\alph{enumi})}
\setcounter{enumi}{2}
\tightlist
\item
  设 N 为 b) 中那样的一个子扩张。依次证明：\(K''\) 是 \(K'\) 的可分代数扩张（因为 L 在 K 上可分）；\(M(x_{1}, \ldots, x_{r})\) 是 \(N(x_{1}, \ldots, x_{r})\) 的可分代数扩张（利用 b)，\(\gamma)\)）；以及 M 是 N 的可分代数扩张。d) 证明 M 是一个纯超越扩张的可分代数扩张（利用 c) 以及 N 是 k 的有限生成可分扩张这一事实）。
\end{enumerate}

